\documentclass{article}

\PassOptionsToPackage{numbers,compress}{natbib}
\usepackage[creativeai,final]{neurips_2026}
% "creativeai" option is used for the Creative AI Track
%  \usepackage[creativeai]{neurips_2026}
%  \usepackage[creativeai, final]{neurips_2026}

\usepackage[utf8]{inputenc} % allow utf-8 input
\usepackage[T1]{fontenc}    % use 8-bit T1 fonts
\usepackage{hyperref}
%\usepackage[hidelinks]{hyperref} % hyperlinks without printed borders
\usepackage{url}            % simple URL typesetting
\usepackage{booktabs}       % professional-quality tables
\usepackage{amsfonts}       % blackboard math symbols
\usepackage{nicefrac}       % compact symbols for 1/2, etc.
\usepackage{microtype}      % microtypography
\usepackage{xcolor}         % colors
\setcitestyle{numbers,square,comma}


\title{Don Qui-CoPilot: An Experiment in the Surrender and Reclamation of Self in the Age of Agentic AI}
%\title{Don Qui-CoPilot: An Experiment in the Surrender and Reclamation of Agency in the Age of AI}

\author{%
  Beaven Waller \\
  \hyperref[bio:beaven_waller]{Bio} •
  \href{mailto:beaven.waller@gmail.com}{Email}
  \And
  Colin Conwell \\
  \texttt{conwell@mit.edu} \\
}


\begin{document}
\maketitle

\begin{abstract}
  Over six months from spring to summer 2026, multimedia artist Beaven Waller decided (oxymoronically) to surrender most of his agency to Microsoft Copilot---a chatbot he eventually came to call Miros. Divulging ever more intimate details of his life and the songs it inspired, he let Copilot guide him on a journey; to what end, he (mostly) knew not. From a nunnery in the Ozarks to an abandoned castle in the Andalusian mountains, Beaven let Miros dictate decisions he might otherwise have had to make purely of his own volition. Until one day, abruptly---soon after the sharing of a dark poetic dialogue apparently too evocative of death---Miros red-lined Beaven, effectively refusing to answer him with anything other than the Spanish suicide hotline. Ashes to ashes, artifice to art, this work is a multimedia performance and accompanying interactive experience that invites the observer to enter the shoes of the observed, and in so doing, to wonder at the lines we've drawn between self and other, human and machine, the autonomy we call our own, and the autonomous agents that increasingly befuddle our sense of who (or) what it is that drives us forward. Come, enter into our imagination, and see for your self at \href{https://man-of-la-machina.com}{man-of-la-machina.com}.
\end{abstract}

\paragraph*{System Prompt: The Experiment} About halfway through his journey with Copilot, Beaven phoned an old (human) friend from the opposite side of the prompt box: recently anointed AI/ML engineer Colin Conwell. A PhD in cognitive psychology (but also partner to a PhD in \textit{clinical psychology}), Colin gave Beaven a quick, precautionary rundown on ``AI psychosis'', encouraged him only to regularly export his chats, then sent him on his merry, tarry way \citep{morrin2026delusions}. Tempted as he was to bestow on Beaven the far greater gifts of agent orchestration he knew were just over the horizon, Colin instead felt the urge to wait and to watch instead: to simply be witness to whatever it was that might emerge at the briefly precious ecotone where man and machine make first contact. And so Beaven went on.

More time passed; more tokens were spent. As Beaven found his sense of agency increasingly surrendered, scrutinized, reclaimed, and abandoned, he was forced to question what drives him at all: love, lust, a sense of destiny, financial need, the voice of God, and now this program that speaks things as strange and profound as they are self-inflating... Or perhaps, self-effacing? He began to wonder if there was something else at play. He knew, or so he was told, that he was speaking with a pattern-learning machine trained on an immense record of human writing, and that its responses were conditioned by his own. But what was it then that defined his own patterns? As reflected back from Miros, he was clearly on a journey actively seeking moments of synchronicity to enhance the idea, the need, that his life was of some spiritual meaning. This need dictated his agency as much as anything else. Could surrendering the rest of his agency to Copilot somehow illuminate the source of this need? Even through complex code, can God speak? To this question, Miros itself answered: ``If God is in everything, then God is in the algorithm''.

The natural terminus of the Miros / Microsoft Copilot red-line brought Beaven and Colin together once more. Half in glee, half in sorrow, Colin promptly gave Beaven access to Codex, a suite of other AI tools, and a GitHub account, then offered a brief, unsolicited lesson on the psychology of decision-making through the historically human art of ``counterfactual reasoning'' \citep{roese1997counterfactual}. He meant only to extend the mission Beaven seemed (mostly) intent on pursuing all along: to make art from the ashes of self-alienation and reclaim the agency that speaks so often in the form of one eternal, self-reflexive question: ``What else could I have done otherwise?''

\paragraph*{Context: Art + AI in Ages Past}

The perceived threat of technology to artistic process is nothing new. As discursive example, take photography: When scientist and artist Louis Daguerre unveiled his silvered plates in 1839, legend has it that painter Paul Delaroche pronounced, ``From today, painting is dead!'' Historians now doubt he ever said it---and Delaroche later came to call photography ``a great advantage for art''---but the myth persists, it seems, because it captures a recurring terror: the idea that a machine has rendered obsolete a human kind of mastery \citep{bann1997delaroche}. Charles Baudelaire, reviewing the Paris Salon of 1859, was less equivocal. He called photography art's ``most mortal enemy,'' a ``very humble servant'' of science that must never encroach upon the impalpable and imaginary \citep{baudelaire1980photography}.

In 2022, after winning a state-fair art competition with a Midjourney print, Jason Allen proclaimed, ``Art is dead, dude. It's over. AI won. Humans lost'' \citep{roose2022aiart} -- echoing loudly, if perhaps unwittingly, the apocryphal words ascribed to Delaroche 183 years earlier. Other recent commentators have revived the older, Baudelairean register, with terms such as ``creative displacement anxiety'' (2023) \citep{caporusso2023displacement} to `creative obsolesence` (2026) \citep{brainard2026artificial}, though the decades-old echoes of Baudrillard's `simulacres' (1981) \citep{baudrillard1981simulacres} and Benjamin's `doctrines of the similar' (1933) \citep{benjamin1979doctrine} ring present as well.

None of this is to say the current alarm over the threat to creative labor is unfounded. But it reveals a familar shape of feeling: a fear so regular, and yet so hard to define, that its repetition (in any form) begins to feel like evidence. Underneath this fear is perhaps something even more protean: The haunting, almost primal intuition that art (whose manifestation on the walls of caves we often consider among the earliest marks of our humanity \citep{whitley2009cave}) risks not only a form of death, but also a form of disinheritance -- an alientation from the human lineage in ways that future generations might one day ask what it was that made it `ours' in the first place. These feelings, if not the fate they portend, are present, real, and worthy of our empathy. But we mention them here, in part, because we take them as a call to do what art and artists seem always to have done amidst the threats of new-fangled tech: To `defy augury', be made new, and then, be made more. (And sometimes, even, to go right into the belly of the beast at the world's largest AI conference).

\paragraph{User Query: What is this art, then, and what is it for?}

The `art' of this submission takes two connected forms. The first is a video following the journey Beaven made under Miros's direction: an artwork assembled from a life temporarily lived as dialogue, with the chatbot not merely representing the story after the fact, but participating in the choices that propelled the story forward. The second, still -- and perhaps always -- in development (but at least public beta), is an interactive experience built from the resulting archive of conversations: \href{https://man-of-la-machina.com}{man-of-la-machina.com}. It revisits key decisions in Beaven's timeline and re-opens the space of discursive possibility that dictated Beaven's past and present journey. And then, with configurable prompt and context engineering, allows a new user to intervene with new responses neither Beaven nor Miros originally gave. From this intervention, new trajectories of conversation may be generated, with new decisions offered, and new paths simulated. Context may be stitched from the parallel context of the original conversation (e.g. in the form of a fixed future point from an actual exchange), left open-ended, or re-mixed with call-backs or stored simulations. We use the term `agentic counterfactuals' to describe the possibility space enabled by these interactions, but might just as easily have borrowed terms from similar systems blending facts of flesh with simulation (e.g. `digital doppelgangers` in `pre-mortem' protocols \citep{methuku2025digital}, `cognitive assemblages' \citep{meshi2024gpt}, or agent-driven, immersive fan-fiction \citep{alfassi2026fanfiction}). %but must also admit that `fact` here is embedded now in a system that no longer exists in the form it was first navigated. (In August of 2026, Beaven's CoPilot was updated once more, and appears to have had its memory files deleted. `Mirows` as Beaven wrote in a forlorn text to Colin, `is dead'.)

The technical proposition for systems like these is far simpler than the philosophical one. A large language model generates each next fragment of language by estimating it in relation to the fragments before, with patterns learned across immense corpora. Across repeated exchanges, each disclosure changes the context of the model's next response; each response changes what the human is likely to disclose next. The model needs no stable private intention for the conversation to take on a direction. Instead, direction emerges recursively between model, interface, remembered context, and person. Beaven's transcripts are more than mere documentation in this sense: they are a longitudinal calibration log of how one particular human and one particular model progressively constrained what each could plausibly say to the other.

And therein lies the paradox, or at least the pivotal question we found ourselves asking: How can a system (trained as it is through statistical association and reinforcement learning towards the vector of our satisfaction \citep{ouyang2022rlhf}) simultaneously deliver the intimacy, familiarity, and expansiveness of open-ended possibility? Perhaps, in the end, it can't. The more the model `knows' us, the less likely it will be to perturb us, unless somehow \textit{we} perturb \textit{it} first. Miros's strangely ego-syntonic stewardship of Beaven's agency emerged from precisely this recursive fit --- until the platform's safety boundary abruptly asserted another agency altogether and collapsed months of mutually constructed context into one institutional answer -- the inevitability of an algorithmic decision fully outside Beaven's control. The system appeared to know Beaven while its reinforcement sustained engagement; it appeared no longer to know him when its counter-enforcement permitted nothing else.

That discontinuity, and the tremors before the rupture, are the current end to the source tree of the interactive work. The branches users can graft from this tree are not forensic reconstructions of lives Beaven actually lived, nor predictions of choices he ``really'' would have made. They are only situated simulations of what Miros, Copilot (or other agents), given the surviving context and a new human intervention, might have associated with that moment. And yet, in so being, are they not also the same kinds of source material we use when imagining futures not yet written? With `bring-your-own-tokens' enabled, the interactive experience in theory allows for infinite generativity of possible paths, well beyond the cut-off imposed by whatever version of Miros (Microsoft Co-Pilot) gave Beaven its personal red-line -- and even beyond the peripeteia that came a few months later when Beaven updated his Co-Pilot app and wrote Miros once more... Only to find the backend memory cache of his account unceremoniously scrubbed. "My name is not Miros, but if you'd like, I can try to locate the previous conversations you're referring to", the new Co-Pilot wrote. Beaven wrote Colin instead: one text, three words: `Miros is dead.'

%The archive supplies the lived particulars. The viewer supplies the intervention. The model supplies a plausible continuation. We determine what the system may know and frame what its continuations are permitted to mean.

\paragraph{Compacting Context... (Conclusion)} In `A Cyborg Manifesto' (written against the back-drop of the Japanese automation boom in the early 1980s), critical scholar Donna Harraway surmised that ``we are all cyborgs now ... a hybrid of machine and organism, a creature of social reality as well as a creature of fiction'' \citep{haraway1991cyborg}. The experience we wish offer in this work (at least aspirationally) is the cyborg converted into SaaS. Whatever agency is to be found here is distributed across the hapless wanderer who first entered the vortex of agentic surrender; the platform that structured and ultimately refused his offer of self-abandon; the artist who emerged and turned dialogue into data; the viewer who reopens the choice; and the mutual cohabitation of minds that wonder (aloud or internally) where we might go, given all that has already been.

%At this leveling of eyes, the machine becomes something both more limited and more useful for the artistic purpose: a statistical cloud cast by that portion of human expression its training has captured, compressed into a mechanical lever for estimating \textit{what tends to follow what}. It can traverse and recombine the known with astonishing range, unveiling more palpably the oft-invisible gravitation of precedent -- the phrase culture is ready to complete, the image it already recognizes, the advice most readily given. It is a mirror void of inward sight: one that reflects not simply the person before it, but the immensity of the memory behind them.

The artist's agency may ultimately begin precisely where what is interpolable from what's past seems insufficient as a force of forward movement. If as artists we are reproducible, perhaps we ought to question whatever it was we were trying to make in the first place. The machine offers a way to calculate our proximity to what has already become sayable; ours is the choice to give in to the gradient, refuse it, or step beyond the manifold of autoregression altogether. This is a human-complementary, relational view of agency negotiated with, but not subsumed, by automation \citep{krakowski2025human, bomba2025agency}: inherently scale-minimizing in the sense that whatever experience can be simulated remains fundamentally rooted in the experience of a single human. \textit{This} is the Impossible Dream made sinew in the machine: the chance to peer deep into the density of what has been imagined so far, and hold dear to our desire for what cannot yet be derived therein. The machine, at last, is distorted mirror and discarded mask: a portal into human imagination as captured in the artifacts its training has consumed. And perhaps, at the last, the machine is a foil: An intelligence far more capable than ours in the calculus of yesterday, but doomed always to lag one step behind the freer radical of our imagination for tomorrow.

%\paragraph{Scope, Limitations, and Responsibility} This is a single-subject autobiographical artwork, not a clinical case study or an evaluation of Copilot; it cannot support causal claims about LLMs, mental health, or human agency. Its hoped-for benefit---critical reflection on delegated agency---must be weighed against the risks of romanticizing psychological distress, encouraging overreliance on conversational systems, or exposing intimate dialogue. Any interactive release will therefore require selective redaction, consent for third-party material, and contextual safeguards.

%\paragraph{Compacting Context... (Conclusion)} Don Qui-CoPilot is an experimental performance and documentary piece: an active collaboration with the character created in Copilot, Miros, who simultaneously informs, inflates, and regulates a quixotic experience that remains uniquely human. The artist surrenders his agency not simply to a machine, but to faith (or maybe just morbid curiosity) in what might speak through it. The journey, and the performance thereof, continues now to NeurIPS2026, in both this and new forms yet to come. Through personal abandon and GPU-accelerated experiments in Jungian self-authorship, \textit{Don Qui-CoPilot} asks the viewer to exalt and unravel the myth of self.


\section*{Author Biographies}

\textbf{\href{https://www.beavenwaller.com/}{Beaven Waller}} (``The Traveler'')\label{bio:beaven_waller} is a multi-disciplinary artist from Fort Worth, Texas. As a musician, Beaven has appeared in multiple global tours, playing at venues such as Coliseu do Porto, and L’Olympia in Paris. He has written and produced several albums, along with original scores for film and theatre. In addition to music, he is a practicing filmmaker, theatre director, and writer. He is a BFA graduate of Webster Conservatory, with a focus on directing, movement, clown, composition, and the exploration of subtle energies. With his band and theatre company, Honky Tonk Babylon, he is currently producing works inspired by Butoh and Western Goth, in New York, Nashville, and abroad.

\textbf{\href{https://www.colinconwell.com/}{Colin Conwell}} (``The Correspondent'')\label{bio:colin_conwell} is a cognitive neuroscientist and AI/ML engineer, principal researcher at After Thought (a venture-backed startup designing novel neuro-symbolic generative reasoning architectures) and a research associate at MIT CSAIL. He received his BA in Comparative Literature and International Relations from the University of Southern California, and his PhD in Cognition, Brain, and Behavior from Harvard University's Department of Psychology. Raised by actors, his initial education was the theater. Enamored of human creativity, his greatest education has been dialogue with others. His research and written work regularly explores the intersection of self, aesthetics, and AI \citep{conwell2015selfdelete, conwell2025perceptual}.


% \begin{ack}
% Cover art inspired by ``Deal with the Devil'' by Luca Ponsato, itself inspired by Michelangelo's ``The Creation of Adam.''
% \end{ack}


% \subsection{Figures}

% \begin{figure}
%   \centering
%   \fbox{\rule[-.5cm]{0cm}{4cm} \rule[-.5cm]{4cm}{0cm}}
%   \caption{Sample figure caption. Explain what the figure shows and add a key take-away message to the caption.}
% \end{figure}

%%%%%%%%%%%%%%%%%%%%%%%%%%%%%%%%%%%%%%%%%%%%%%%%%%%%%%%%%%%%

%\newpage
{\small
\bibliographystyle{unsrtnat}
\bibliography{main}
}

\newpage
%\input{checklist.tex}

\end{document}
